\section{Closed-form gradient compliance}\label{sec:compliance}

\subsection{Inverse and conditions}\label{sec:compliance-inverse}

By~\eqref{eq:block-rules}, the compliances of the cubic tensor~\eqref{eq:cub-tensor} and
of the hexagonal tensor~\eqref{eq:hex-tensor} are
\begin{equation}
\begin{aligned}
\tD^{-1}&=\sum_{a,b}(M^{-1})_{ab}\tU{T_1}{a}{b}+\sum_{a,b}(N^{-1})_{ab}\tU{T_2}{a}{b}
+\frac{\tP{E}}{m_E}+\frac{\tP{A_2}}{m_A},\\
\tD^{-1}&=\sum_{a,b}(\hexA^{-1})_{ab}\tU{A_2}{a}{b}+\sum_{a,b}(\hexF^{-1})_{ab}\tU{E_1}{a}{b}
+\sum_{a,b}(\hexG^{-1})_{ab}\tU{E_2}{a}{b}
+\frac{\tP{A_1}}{m_{A_1}}+\frac{\tP{B_1}}{m_{B_1}}+\frac{\tP{B_2}}{m_{B_2}}.
\end{aligned}\label{eq:compliance}
\end{equation}
The blocks have size at most three (cubic) and four (hexagonal), so their inverses are
explicit by cofactors, and the units are explicit in components (supplementary
material, section~S1); the same formulas invert the non-symmetric elements of the
algebras, whose blocks are not symmetric. A cubic tensor is invertible if and only if
$\det M\,\det N\,m_Em_A\ne0$, and positive definite if and only if, by Sylvester's
criterion on the blocks,
\begin{equation}
M_{hh}>0,\quad M_{hh}M_{ss}-M_{hs}^2>0,\quad \det M>0,\quad N_{33}>0,\quad \det N>0,
\quad m_E>0,\quad m_A>0.\label{eq:cub-pd}
\end{equation}
These are seven conditions, as many as Lazar et al.\ obtain from Sylvester's criterion
on the $5\times5$ and $3\times3$ diagonal blocks of their matrix~\cite[appendix~A,
eqs.~(A.7)]{Lazar2022}, but each involves a single irreducible block and has degree at
most three in the moduli. A hexagonal tensor is invertible if and only if $\det\hexA\,\det\hexF\,\det\hexG\,
m_{A_1}m_{B_1}m_{B_2}\ne0$ and positive definite if and only if $\hexA$, $\hexF$ and
$\hexG$ are positive definite and $m_{A_1}$, $m_{B_1}$, $m_{B_2}$ are positive; in the
components of Auffray et al., \eqref{eq:hex-auffray} turns these into the positive
definiteness of $A^{(11)}$, $A^{(11)}+\eta_{\mathrm A}A_c$, $H^{(6)}+f(J^{(4)})$ and $J^{(4)}$.
By~\eqref{eq:hex-extra},
\begin{equation}
\frac{\tP{B_1}}{m_{B_1}}+\frac{\tP{B_2}}{m_{B_2}}
=\frac12\Big(\frac1{m_{B_1}}+\frac1{m_{B_2}}\Big)(\tP{B_1}+\tP{B_2})
+\frac18\Big(\frac1{m_{B_1}}-\frac1{m_{B_2}}\Big)\tHsix,\label{eq:hex-extra-compliance}
\end{equation}
so the extra modulus of $\tD^{-1}$ is $-\Dhex/(m_{B_1}m_{B_2})$: the compliance of a
transversely isotropic tensor is transversely isotropic, and that of an isotropic tensor
is the isotropic inverse~\eqref{eq:iso-tensor}.

\subsection{Back to the constants of Lazar et al.}\label{sec:compliance-lazar}

The map~\eqref{eq:cub-lazar} is inverted in closed form. With
$u=M_{hh}-m_A$, $x=N_{33}-m_A$, $y=N_{22}-m_E$ and
$X=M_{ss}-\sqrt6M_{hs}+\tfrac32M_{hh}-\tfrac52N_{33}$, $Y=M_{rr}-N_{22}$,
$Z=M_{sr}-\tfrac{\sqrt6}{2}M_{hr}-\tfrac{\sqrt{10}}{2}N_{32}$,
\begin{equation}
\begin{gathered}
a_1=\tfrac1{15}\big(2X-5Y+\sqrt5Z\big),\quad
a_2=\tfrac1{30}\big(X+5Y+2\sqrt5Z\big),\quad
a_3=\tfrac1{30}\big(4X+5Y-4\sqrt5Z\big),\\
a_4=\tfrac16(m_A+2m_E),\quad a_5=\tfrac13(m_A-m_E),\quad
a_6=\tfrac16\big(2x+y-2\sqrt2N_{32}\big),\\
a_7=\tfrac16\big(2x-2y+\sqrt2N_{32}\big),\quad
a_8=\tfrac13\big(x+2y+2\sqrt2N_{32}\big),\quad
a_{11}=\tfrac12(5u-9x),\\
a_9=\tfrac12(x-u)+\tfrac16\big(\sqrt6M_{hs}+\sqrt{30}M_{hr}+3\sqrt2N_{32}\big),\\
a_{10}=\tfrac12(x-u)+\tfrac1{12}\big(2\sqrt6M_{hs}-\sqrt{30}M_{hr}-3\sqrt2N_{32}\big).
\end{gathered}\label{eq:cub-lazar-inverse}
\end{equation}
For an isotropic tensor $u=x=y=0$, $X=M_{ss}-m_3$, $Y=M_{rr}-m_2$ and $Z=M_{sr}$, and
the first line, with $a_4$ and $a_5$, inverts~\eqref{eq:iso-lazar}. The conditions
\eqref{eq:cub-pd} in the constants $a_i$ follow by inserting~\eqref{eq:cub-lazar}.

\subsection{Atomistic data}\label{sec:compliance-data}

Table~\ref{tab:block-moduli} gives the block moduli of five cubic data sets: the MEAM
constants of Al, Cu, Fe and W of Lazar et al.~\cite[table~1]{Lazar2022}, converted to
$a_1,\dots,a_{11}$ by their eq.~(60), and the EAM constants of Cu of Po et
al.~\cite[table~1]{Po2019}. Po et al.\ print the Al and Cu MEAM constants too, with five
significant figures instead of five decimals; for the entries below $0.1$ in magnitude
their digits are finer, and we take each entry from the finer of the two printings. Shodja et al.~\cite{Shodja2018} evaluated cubic constants by
density functional theory under a central-force assumption, which imposes equalities
among them; we
use the complete tensors. The data of Al and of both potentials for Cu are positive
definite. Those of Fe and W are not: Lazar et al.\ find one negative eigenvalue, threefold, a
root of the quadratic factor~\cite[eqs.~(130)--(133) and table~5]{Lazar2022}. In the basis
that factor is the $\irr{T_2}$ block, the only one that fails, with $\det N=\NumFeDetNT$ and $\NumWDetNT$~(eV/\AA)$^2$. Its
entry $N_{33}$, the $\irr{T_2}$ part of $\Hrm3$, is negative ($\NumFeNthreethreeT$ and
$\NumWNthreethreeT$~eV/\AA), but for W it nearly vanishes, and the indefiniteness comes
from the coupling $N_{32}=\NumWNthreetwoT$~eV/\AA. With~\eqref{eq:cub-lazar}, the two
conditions $N_{33}>0$ and $\det N>0$ are explicit in the constants $a_i$, and a fit of a
potential could impose them. Indefinite gradient tensors
are common among interatomic potentials~\cite{Admal2017,Lazar2022}; Admal et
al.~\cite[eq.~(3.11)]{Admal2017} propose instead the positive definiteness of
$\tfrac15R^2\,\tC\ot\kron+\tD$ for a sphere of radius $R$. Both terms are cubic, so the
condition is a generalised eigenvalue problem per block; it holds for
$R>\NumFeNeighborhoodRadius$~\AA\ (Fe) and $R>\NumWNeighborhoodRadius$~\AA\ (W), that is
$\NumFeNeighborhoodRatio\,b$ and $\NumWNeighborhoodRatio\,b$ with $b$ the
nearest-neighbour distance, consistent with their finding that it holds beyond $1.2\,b$ for
all the potentials they examine, Fe and W included~\cite[section~3]{Admal2017}. The Cu MEAM tensor
is nearly singular, with condition number $\NumCuMEAMCond$, because its $m_E$ is close to
zero; the rounding of the published constants fixes $m_E$ only to within
$\pm\NumCuMEAMmERoundingPercent\,\%$. These are properties of the potentials, not of the
method. The compliance blocks of the three positive definite sets are in the
supplementary material (table~S3).

\subsection{Verification}\label{sec:compliance-verification}

The results were checked by routes other than those that produced them. Blocks, inverses
and averages are compared with the same operations on the $18\times18$ matrix of the
tensor restricted to $\Vsg$; the multiplication tables, the inverses of general elements
and the closed averages are also verified in exact arithmetic; the formulas printed in
sections~\ref{sec:setting}--\ref{sec:averages} and appendix~\ref{app:lazar} are restated
in the tests as printed and checked against the code. Published results enter after reconstruction: the constants
and eigenvalues of Lazar et al.~\cite{Lazar2022}, the matrices of Auffray
et al.~\cite{Auffray2013} and the tensors of Monchiet and Bonnet~\cite{Monchiet2011a,Monchiet2013}
are reproduced (except the eigenvalue $\lambda_7^D$ of Cu in~\cite[table~5]{Lazar2022},
printed as a repetition of $\lambda_4^D$; the $\irr{T_1}$ block gives
$\NumCuMEAMLambdaSeven$~eV/\AA), and where a printed formula does not hold, both the printed and the
corrected form are checked. Three such formulas enter this paper: the isotropic inverse of
Monchiet and Bonnet~\cite[eq.~(2.13)]{Monchiet2011a}, whose terms in their $J_3$ and $J_4$
need the opposite signs; the average over the sphere of the sixth-order Eshelby tensor of
Ma et al.~\cite[eq.~(3.18)]{Ma2018} (appendix~\ref{app:eshelby}); and, in the scheme of
Solyaev~\cite{Solyaev2023}, the deviatoric term of the fourth-order average in his eq.~(20),
half that of Ma et al., his eq.~(21), which inherits their eq.~(3.18), and his eq.~(22), printed without the
squares (section~\ref{sec:application}). We distinguish three levels: verification, against exact
mathematics or an independent route of our own; cross-comparison, against a
reconstructed published result; and consistency, through limit cases, invariances and
inequalities on samples. No comparison is with measurements: the inputs are published
atomistic constants and, in the hexagonal class, tensors of granular micromechanics. The tables, figures and
computed numbers are written by scripts of the public code (Data accessibility), whose
tests compare each file with its generator; section~S2 of the supplementary material
maps each claim to its key tests, their reference and their level.

\begin{table}[pos=tbp]
\caption{Block moduli~\eqref{eq:cub-tensor} of the five cubic data sets (eV/\AA), from
the MEAM constants of Lazar et al.~\cite[table~1]{Lazar2022} and the EAM constants of
Po et al.~\cite[table~1]{Po2019} by~\eqref{eq:cub-lazar}; for Al and Cu MEAM each entry
is taken from the finer of the two printings. Last rows: positive
definiteness, with the failing block, and condition number of $\tD$ (largest eigenvalue over
smallest eigenvalue modulus). The last digit shown is uncertain by at most five units
given the precision of the published constants; $\dagger$ marks the entries whose
relative uncertainty exceeds 1\,\%, here the $m_E$ of Cu MEAM ($\pm\NumCuMEAMmERoundingPercent\,\%$)
and the condition number it sets. The Cu EAM constants obey the relations
$D_{2,4}=D_{3,5}$ and $D_{2,5}=D_{16,17}$ that Admal et al.~\cite{Admal2017} find for
pair functionals, that is, $a_3=4a_2$ and $a_1=a_5$.}\label{tab:block-moduli}
\centering\small
\input{block_moduli}
\end{table}
